% Generated draft: requires mathematical and source review.
\documentclass[11pt]{article}
\usepackage[margin=25mm]{geometry}
\usepackage{fontspec}
\setmainfont{Latin Modern Roman}
\usepackage{amsmath,amssymb,mathtools,unicode-math}
\setmathfont{Latin Modern Math}
\usepackage{enumitem,needspace,xurl,hyperref}
\setlength{\parindent}{0pt}
\setlength{\parskip}{5pt}
\setlength{\emergencystretch}{3em}
\newcommand{\N}{\mathbb{N}}
\newcommand{\Z}{\mathbb{Z}}
\newcommand{\Q}{\mathbb{Q}}
\newcommand{\R}{\mathbb{R}}
\newcommand{\C}{\mathbb{C}}
\newcommand{\F}{\mathbb{F}}
\newcommand{\A}{\mathbb{A}}
\newcommand{\PP}{\mathbb{P}}
\newcommand{\E}{\mathbb{E}}
\newcommand{\eps}{\varepsilon}
\newcommand{\dd}{\,\mathrm{d}}
\begin{document}
\section*{UnsolvedMath ID: 701281}
\textbf{Source problem code:} catalogue entry
\subsection{Definitions and mathematical statement}
Let \(X\) be an infinite set, let \(F=F(X)\) be the free group with free basis \(X\), and let \(G=\operatorname{Aut}(F)\) be its full automorphism group. For \(g\in G\), define its normal closure in \(G\) by
\[
\langle\!\langle g\rangle\!\rangle_G
=\bigl\langle hgh^{-1}:h\in G\bigr\rangle,
\]
the smallest normal subgroup of \(G\) containing \(g\). An involution is an element of order exactly two.

For every infinitely generated free group \(F\), the problem asks the following two questions:
\begin{enumerate}
\item[(a)] Does there exist \(g\in\operatorname{Aut}(F)\) such that
\[
\operatorname{Aut}(F)=\langle\!\langle g\rangle\!\rangle_{\operatorname{Aut}(F)}?
\]
\item[(b)] Does there exist \(\tau\in\operatorname{Aut}(F)\) such that
\[
\tau\ne 1,\qquad \tau^2=1,\qquad
\operatorname{Aut}(F)=\langle\!\langle\tau\rangle\!\rangle_{\operatorname{Aut}(F)}?
\]
\end{enumerate}
The source specifies no restriction on the infinite cardinality of \(X\). It does not require the involution in part (b) to permute a free basis.

As background, the supplied record attributes to M. Bridson and K. Vogtmann the assertion that, for a finite-rank free group \(F_n\), its automorphism group is normally generated by an involution that permutes some free basis. The record describes this result as ``recently'' obtained; this is source wording, not a present-day literature assessment. It also reports the analogous existence of such involutions for automorphism groups of infinitely generated free nilpotent groups. These background assertions are retained as supplied reports: the record gives neither separate references nor precise parameter restrictions for them, including the finite-rank range and the nilpotency class. They are not additional hypotheses on \(F\).
\subsection{Short English statement}
For an infinitely generated free group, is its full automorphism group normally generated by one element? Can that element always be chosen to have order two? The supplied registry labels the problem open; this label and the accompanying background results have not been independently verified.
\subsection{Sources}
\begin{itemize}
\item Supplied local UnsolvedMath registry record, permanent ID \(701281\), exact source problem code \texttt{catalogue entry}, titled \emph{Kourovka Notebook Problem 16.90}. The record labels the problem \emph{open}; its supplied creation and update dates are 2026-10-16.
\item \emph{The Kourovka Notebook (Unsolved Problems in Group Theory)}, Problem 16.90, 16th issue (2006); the supplied record also identifies the 21st-issue edition (2026), arXiv:1401.0300v48. The collection metadata names E. I. Khukhro and V. D. Mazurov as editors. Proposed by V. Tolstykh in 2006.
\item The supplied statement attributes its finite-rank background result to M. Bridson and K. Vogtmann and reports a corresponding result for infinitely generated free nilpotent groups without naming its authors. No separate bibliographic references or external links for these background assertions are supplied.
\end{itemize}
\end{document}
